\chapter{Method}

The findings in this manual were arrived at through a deliberately
restrictive method. The method is worth stating explicitly because it
shapes which claims appear, and because the same discipline is what
allows the reader to trust those claims.

\section{Claim Discipline}

Three categories. They are never blurred.

\textbf{Spec.} A claim backed by Apple's public reference or by another
durable specification. Spec claims are useful but not sufficient. The
reference describes intent; what is on disk is what the running system
writes, and the two occasionally disagree.

\textbf{Observation.} A claim backed either by a controlled experiment
on a real APFS image, or by behaviour that converges across multiple
independent open-source readers. Convergence across six readers
(\texttt{linux-apfs-rw}, \texttt{apfs-fuse}, \texttt{dissect.apfs},
\texttt{go-apfs}, \texttt{libfsapfs}, the Sleuth Kit's APFS module)
plus the reference is harder to fault than the reference alone.

\textbf{Hypothesis.} A claim that is plausible and useful for design
but has no probe or convergence behind it yet. Hypothesis does not
license product behaviour. If the parser would emit a row only on a
Hypothesis-class claim, the parser does not emit the row.

The two rules that follow from these labels.

\paragraph{Where Spec and Observation disagree, Observation wins.}
Section~\ref{ssec:xfield-spec-vs-impl} on the xfield cursor rule is one
example; section~\ref{ssec:fs-tree-internal} on the FS-tree internal
edge is another. In both cases the reference is easy to read in a
direction that the converged code does not support. The converged
code is right.

\paragraph{Where a Hypothesis matters, the smallest decisive probe is
named.} The manual does not leave open questions floating; it leaves
them attached to the experiment that would close them.

\section{Oracle Plurality}

There is no single oracle for an APFS indexer. The phrase "the
oracle" is a tell that someone has glued mismatched validation sources
together and called the result authoritative. Each feature owns its
own oracle.

\begin{description}
  \item[Namespace.] POSIX traversal (\texttt{os.walk}) of the chosen
    volume at the chosen state. Path bytes, entry type, hard-link
    grouping by inode, symlink target.
  \item[Logical size.] \texttt{st\_size} per inode; for symlinks, the
    target byte length returned by \texttt{readlink}. This is the
    public, stable per-file size that user-space tools agree on.
  \item[FS-record body.] Cross-tool field equality between
    independent decoders against the same raw container. The
    \texttt{go-apfs identitydump} tool plus an independent Python
    decoder against the same \texttt{/dev/rdiskN} is the form the
    project uses.
  \item[Allocated size.] \texttt{st\_blocks * 512}, and only on cases
    that have an explicit validated mapping from raw fields to it.
  \item[Incremental correctness.] A fresh full reparse of the same
    selected state. Any cache reuse must equal this.
  \item[Boot-root semantics.] The user-visible macOS namespace ---
    System/Data join, firmlinks, current snapshot --- and only when
    the product mode targets it.
\end{description}

A probe that does not name its oracle does not produce evidence.

\section{Paired Probes}

The most informative experiments build the fixture, capture every
oracle, run the candidate parser, and compare --- all in one execution
of one script. This rules out the kind of subtle drift that breaks an
experiment: a fixture rebuilt with slightly different operations, a
volume remounted between probes, a parser pointed at a different
container than the oracle.

Two patterns in particular keep recurring.

\paragraph{Detach-and-reattach.} Build the fixture mounted; capture
the mounted POSIX oracle; \texttt{hdiutil detach}; \texttt{hdiutil
attach -nomount -readonly}; run the raw parser; diff. The detach
fixes the visible checkpoint, so the raw parser and the mounted
oracle agree on what state to compare.

\paragraph{Cross-tool agreement.} Run the candidate parser and an
independent third-party reader against the same raw container in the
same probe. Record the parser's resolved \texttt{(oid, paddr, xid)}
samples; replay the object-header validation rules at every resolved
\texttt{paddr}; confirm the two tools agree on the FS-tree root OID.
A divergence in \emph{which} checkpoint each tool selected is itself
recorded as a caveat, not silenced.

\section{The Smallest Useful Fixture}

Probes prefer the smallest fixture that exercises the question, not
the largest fixture that might. The proof fixture used through most
of this manual is intentionally tiny: two directories, a renamed
file, a moved file, a hard link, a sparse 1 MiB file, a clone, an
append after the clone, and a symlink. Roughly seven user-visible
entries.

That fixture is enough to surface:

\begin{itemize}
  \item the FS-tree internal-OID-vs-paddr distinction (when the tree
    grows past one leaf node),
  \item the xfield cursor rule (the sparse inode has three xfields, one
    of which has a non-multiple-of-eight \texttt{x\_size}),
  \item the logical-size precedence over ordinary, sparse, clone,
    hard-linked, and symlink rows,
  \item the unique-inode aggregate policy (hard links must not
    double-count).
\end{itemize}

Larger fixtures appear only when a specific question requires a feature
the small fixture cannot express, such as the case-insensitive vs
case-sensitive volume comparison in chapter~9 or the compressed-file
precedence test in chapter~8.

\section{Fail Closed By Default}

A parser that emits rows for cases it does not understand is worse
than one that emits nothing. The user can recognise a missing row and
look elsewhere; a wrong row propagates into a directory aggregate,
into a treemap, and into the decisions the user makes from the
treemap. Fail-closed is not an aesthetic preference. It is the only
way to bound the cost of being wrong.

Every step in the parser is wrapped in a typed validation gate. The
single chokepoint for object-header validation is one function
(\texttt{validate\_object\_block} in the Rust crate); every reader
calls it before trusting any block. The single chokepoint for
record-body validation is one function (\texttt{decode\_fs\_record});
every emission path calls it before trusting any record. Each gate
returns a typed error with the rule that failed.

Outside the allowlist, the parser does not partially succeed. It does
not "try its best." It returns a reason and the fallback path
(chapter~11) takes over.

\section{The Style Of A Finding}

The manual presents a finding the same way each time: a one-line
statement, the evidence, the implementation consequence, and the
limits of what was checked. The structure is borrowed from how working
proofs are written: the claim must survive standing alone, but the
proof must be present.
